\chapter{RESULTS AND DISCUSSION}

\section{Output Screenshots with Explanations}
The developed full-stack web dashboard and Three.js 3D digital twin provide real-time visual feedback and supervisory control across six dedicated modules, illustrated in Figures \ref{fig:ui_views_1} through \ref{fig:ui_calib}:
\begin{itemize}
    \item \textbf{Interactive 3D Table Twin \& Click-to-Aim (Figures \ref{fig:ui_3d}, \ref{fig:ui_laser}):} Renders the regulation $1168.4 \times 2341.88\,\text{mm}$ baize surface, six drop pockets, baulk line, 'D' zone, and live 3D ball meshes synchronized with YOLOv8s coordinates. Clicking any cloth point projects a virtual ray from the top-right gantry mast and dispatches spline-interpolated pan/tilt angles to the ESP32.
    \item \textbf{Planar Homography \& Hardware Controls (Figures \ref{fig:ui_homography}, \ref{fig:ui_servo}):} Enables interactive 4-point corner pocket calibration to compute $\mathbf{H}$, alongside direct servo sliders, IP diagnostics (`10.21.46.190`), and laser safety toggles.
    \item \textbf{YOLOv8s Multi-Ball Detection \& Autonomous Referee (Figures \ref{fig:ui_yolo}, \ref{fig:ui_referee}):} Overlays real-time bounding boxes, confidence scores, and millimeter coordinates on overhead frames while the autonomous referee tracks inter-frame motion ($\Delta I_t$), isolates displaced balls after settlement, and provides one-click ``Aim Laser'' restoration buttons.
    \item \textbf{9-Point Spline Calibration Wizard (Figure \ref{fig:ui_calib}):} Guides the operator through fine-tuning the laser spot at nine fiducial nodes to generate the 2D bivariate cubic spline surface.
\end{itemize}

\section{Performance Evaluation}
\subsection{Computer Vision Detection Accuracy and Performance}
The fine-tuned YOLOv8s object detection network was rigorously evaluated on a held-out test dataset comprising 350 multi-ball snooker table images captured under varied ambient and artificial illumination conditions. Table \ref{tab:detection_metrics} presents the per-class performance metrics, including Precision ($P$), Recall ($R$), F1-Score, and Mean Average Precision at an IoU threshold of 0.50 ($\text{mAP}_{0.5}$) and over the range 0.50 to 0.95 ($\text{mAP}_{0.5:0.95}$).

\begin{table}[H]
\centering
\caption{YOLOv8s Ball Detection Performance Metrics Across Classes}
\label{tab:detection_metrics}
\resizebox{\textwidth}{!}{
\begin{tabular}{|l|c|c|c|c|c|}
\hline
\textbf{Ball Class} & \textbf{Precision (\%)} & \textbf{Recall (\%)} & \textbf{F1-Score} & \textbf{$\text{mAP}_{0.5}$ (\%)} & \textbf{$\text{mAP}_{0.5:0.95}$ (\%)} \\
\hline
Orange Ball (Object) & 96.4 & 95.8 & 0.961 & 97.2 & 78.4 \\
White Ball (Cue)     & 97.8 & 96.2 & 0.970 & 98.1 & 81.2 \\
Yellow Ball (2 pts)  & 93.5 & 91.0 & 0.922 & 92.8 & 72.5 \\
Green Ball (3 pts)   & 92.1 & 89.6 & 0.908 & 91.4 & 70.1 \\
Blue Ball (5 pts)    & 94.8 & 93.2 & 0.940 & 95.0 & 75.8 \\
Pink Ball (6 pts)    & 93.0 & 91.5 & 0.922 & 93.2 & 73.0 \\
\hline
\textbf{Overall Mean} & \textbf{94.6} & \textbf{92.9} & \textbf{0.937} & \textbf{94.6} & \textbf{75.2} \\
\hline
\end{tabular}
}
\end{table}

The model achieved an overall $\text{mAP}_{0.5}$ of $94.6\%$. Notably, the Orange object balls and White cue ball exhibited the highest detection accuracy ($97.2\%$ and $98.1\%$ respectively) due to their sharp chrominance separation from the green tournament felt. The Green ball exhibited slightly lower recall ($89.6\%$) due to occasional chromatic blending with the green felt under low-angle indoor illumination.

\section{Comparison with Existing Systems}
To benchmark YOLOv8s against established computer vision methodologies in literature, a comparative study was conducted using the same test dataset. As shown in Table \ref{tab:model_comparison}, YOLOv8s achieves the optimal balance between high precision and low inference latency.

\begin{table}[H]
\centering
\caption{Comparison of Detection Models on Snooker Table Dataset}
\label{tab:model_comparison}
\resizebox{\textwidth}{!}{
\begin{tabular}{|l|c|c|c|c|}
\hline
\textbf{Method / Architecture} & \textbf{Precision (\%)} & \textbf{Recall (\%)} & \textbf{Latency (ms)} & \textbf{Model Size (MB)} \\
\hline
Circular Hough Transform (CHT) \cite{ref4} & 78.2 & 72.4 & 18.5 & --- \\
Faster R-CNN (ResNet-50) \cite{ref33}      & 93.8 & 91.2 & 98.4 & 108.4 \\
Mask R-CNN \cite{ref29}                    & 95.1 & 92.8 & 124.0 & 168.2 \\
U-Net Custom Segmentor \cite{ref1}         & 96.2 & 94.5 & 28.2 & 48.6 \\
\textbf{YOLOv8s (Proposed)}                & \textbf{94.6} & \textbf{92.9} & \textbf{16.8} & \textbf{22.5} \\
\hline
\end{tabular}
}
\end{table}

While U-Net semantic segmentation achieves marginally higher precision ($+1.6\%$), YOLOv8s operates with a significantly smaller parameter footprint ($22.5\,\text{MB}$) and executes at over $59\,\text{FPS}$ ($16.8\,\text{ms}$ latency on GPU), ensuring seamless real-time tracking.

\section{Planar Homography Metric Localization Accuracy}
The geometric fidelity of the computed homography matrix $\mathbf{H}$ was evaluated across nine spatial test locations distributed over the $1168.4 \times 2341.88\,\text{mm}$ table. Physical ground-truth positions were recorded using a high-precision laser distance meter ($\pm 0.5\,\text{mm}$ accuracy).

Table \ref{tab:homography_accuracy} presents the measured spatial localization errors before and after applying the spherical elevation parallax compensation derived in Section 4.2.

\begin{table}[H]
\centering
\caption{Spatial Localization Accuracy Before and After Elevation Parallax Compensation}
\label{tab:homography_accuracy}
\resizebox{\textwidth}{!}{
\begin{tabular}{|l|c|c|c|c|}
\hline
\textbf{Table Test Location} & \textbf{Ground Truth $(X, Y)$ (mm)} & \textbf{Raw Error (mm)} & \textbf{Parallax-Corrected Error (mm)} & \textbf{Improvement} \\
\hline
Top-Right Corner (P1)   & $(50.0, 50.0)$       & $11.8$ & $1.9$ & $83.9\%$ \\
Top-Center Node (P2)    & $(584.2, 50.0)$      & $8.4$  & $1.5$ & $82.1\%$ \\
Top-Left Corner (P3)    & $(1118.4, 50.0)$     & $12.4$ & $2.1$ & $83.1\%$ \\
Mid-Right Pocket (P4)   & $(50.0, 1170.9)$     & $9.6$  & $1.8$ & $81.2\%$ \\
Table Center Node (P5)  & $(584.2, 1170.9)$    & $1.2$  & $1.1$ & $8.3\%$ \\
Mid-Left Pocket (P6)    & $(1118.4, 1170.9)$    & $9.8$  & $1.9$ & $80.6\%$ \\
Bottom-Right Corner (P7)& $(50.0, 2291.9)$     & $11.5$ & $2.2$ & $80.9\%$ \\
Bottom-Center Node (P8) & $(584.2, 2291.9)$    & $8.6$  & $1.7$ & $80.2\%$ \\
Bottom-Left Corner (P9) & $(1118.4, 2291.9)$    & $12.1$ & $2.2$ & $81.8\%$ \\
\hline
\textbf{Average Error}  & ---                  & \textbf{9.49 mm} & \textbf{1.82 mm} & \textbf{80.8\% Reduction} \\
\hline
\end{tabular}
}
\end{table}

Without elevation parallax correction, radial projection caused the apparent ball centers to shift outward by up to $12.4\,\text{mm}$ near table corners. Applying our closed-form parallax correction reduced the mean spatial error from $9.49\,\text{mm}$ to $1.82\,\text{mm}$---well below the $26.25\,\text{mm}$ ball radius, ensuring sub-ball localization accuracy.

\section{Laser Aiming Precision: Theoretical Kinematics vs. Spline Calibration}
To evaluate the targeting accuracy of the steerable laser beam, physical beam contact spots were measured against actual ball targets across the nine calibration nodes.

Table \ref{tab:spline_vs_linear} compares the targeting error of standard theoretical inverse kinematics (Equations \ref{eq:pan_ideal} and \ref{eq:tilt_ideal}) against the 2D Bivariate Cubic Spline interpolation model (Equations \ref{eq:spline_pan} and \ref{eq:spline_tilt}).

\begin{table}[H]
\centering
\caption{Laser Spot Targeting Deviation: Theoretical Kinematics vs. 2D Bivariate Spline}
\label{tab:spline_vs_linear}
\resizebox{\textwidth}{!}{
\begin{tabular}{|l|c|c|c|}
\hline
\textbf{Target Table Location} & \textbf{Theoretical Kinematics Error (mm)} & \textbf{2D Spline Error (mm)} & \textbf{Error Reduction} \\
\hline
Top-Right Home Node (0, 0)       & $4.2$  & $1.2$ & $71.4\%$ \\
Top-Center Node (584, 0)         & $8.6$  & $1.8$ & $79.1\%$ \\
Top-Left Pocket (1168, 0)        & $14.5$ & $2.3$ & $84.1\%$ \\
Mid-Right Pocket (0, 1171)       & $11.2$ & $1.9$ & $83.0\%$ \\
Table Center Node (584, 1171)    & $16.8$ & $2.1$ & $87.5\%$ \\
Mid-Left Pocket (1168, 1171)     & $19.4$ & $2.5$ & $87.1\%$ \\
Bottom-Right Pocket (0, 2342)    & $18.2$ & $2.4$ & $86.8\%$ \\
Bottom-Center Node (584, 2342)   & $21.5$ & $2.6$ & $87.9\%$ \\
Bottom-Left Pocket (1168, 2342)  & $22.6$ & $2.8$ & $87.6\%$ \\
\hline
\textbf{Mean Targeting Error}    & \textbf{15.22 mm} & \textbf{2.18 mm} & \textbf{85.7\% Improvement} \\
\hline
\end{tabular}
}
\end{table}

As the distance from the top-right gimbal mast increases toward the bottom-left pocket ($2616\,\text{mm}$ Euclidean distance), theoretical inverse kinematics exhibits severe targeting drift, culminating in a $22.6\,\text{mm}$ error due to cumulative servo horn flexure, gear train backlash, and gravity droop. In contrast, the 2D Bivariate Cubic Spline maintains an exceptional mean targeting accuracy of $2.18\,\text{mm}$ ($\le \pm 2.8\,\text{mm}$ maximum deviation), satisfying the strict non-functional accuracy requirement (NFR-1).

\section{System Latency and Execution Benchmarking}
To verify real-time performance, end-to-end latency was profiled from the moment post-shot motion cessation is confirmed to the physical projection of the laser spot onto the baize. Table \ref{tab:latency_breakdown} summarizes the execution timing breakdown.

\begin{table}[H]
\centering
\caption{End-to-End Latency Breakdown for Ball Restoration Pipeline}
\label{tab:latency_breakdown}
\resizebox{\textwidth}{!}{
\begin{tabular}{|l|l|c|}
\hline
\textbf{Pipeline Stage} & \textbf{Component / Function} & \textbf{Latency (ms)} \\
\hline
Frame Acquisition       & Camera sensor capture and JPEG stream decoding & $32.4 \pm 4.2$ \\
Deep Learning Inference & YOLOv8s forward pass ($1024 \times 1024$) + NMS & $48.6 \pm 5.1$ \\
Coordinate Rectification& Planar homography inversion + Parallax correction & $3.8 \pm 0.4$ \\
Bipartite State Match   & Hungarian Assignment Algorithm & $5.2 \pm 0.8$ \\
Spline Angle Calculation& 2D Bivariate Spline evaluation for Pan/Tilt & $2.1 \pm 0.2$ \\
Network Transmission    & HTTP POST JSON payload to ESP32 (`10.21.46.190`) & $34.5 \pm 6.8$ \\
Mechanical Actuation    & SG90/MG996R servo slew to target angles ($30^\circ$ swing) & $180.0 \pm 25.0$ \\
\hline
\multicolumn{2}{|r|}{\textbf{Total End-to-End Response Time:}} & \textbf{306.6 $\pm$ 42.5 ms} \\
\hline
\end{tabular}
}
\end{table}

The total execution time of approximately $307\,\text{ms}$ is substantially lower than our design threshold of $1.0\,\text{second}$ (NFR-2), enabling instantaneous, fluid laser guidance for the match referee.

\section{Foul Detection and Restoration Experimental Test Cases}
The autonomous referee engine was validated across four real-world match scenarios:

\subsection{Test Case 1: Standard Miss Foul}
\begin{itemize}
    \item \textbf{Scenario:} Striker takes an escape stroke from the 'D' zone toward an isolated orange ball near the top cushion, missing it completely. The cue ball travels to the bottom-right rail while all object balls remain untouched.
    \item \textbf{System Response:} Optical motion sensor tracks movement until the cue ball rests. State diffing identifies that all orange balls have $C_{ij} < 3.2\,\text{mm}$ (classified as Stationary). The cue ball exhibits displacement $\|\mathbf{d}\| = 1840\,\text{mm}$. The referee calls ``Foul and a Miss''. Clicking ``Restore Cue Ball'' projects the laser directly onto its pre-shot coordinate $(X, Y) = (584.2, 737.0)\,\text{mm}$. The referee replaces the white ball onto the laser spot in under $8\,\text{seconds}$.
\end{itemize}

\subsection{Test Case 2: Multi-Ball Disturbance (Cluster Scatter)}
\begin{itemize}
    \item \textbf{Scenario:} An attempted thin cut results in a heavy collision, scattering three orange balls and deflecting the blue ball. A foul is called.
    \item \textbf{System Response:} Bipartite matching matches post-shot balls to pre-shot coordinates. The system highlights four displaced balls in the UI table (Orange 1: $142\,\text{mm}$, Orange 2: $88\,\text{mm}$, Orange 3: $310\,\text{mm}$, Blue: $54\,\text{mm}$). The referee clicks through the restoration list; the laser sequentially steps to each pre-shot coordinate, allowing orderly, dispute-free ball replacement.
\end{itemize}

\subsection{Test Case 3: Accidental Ball Pocketing}
\begin{itemize}
    \item \textbf{Scenario:} Cue ball accidentally drops into the bottom-left corner pocket during an escape attempt.
    \item \textbf{System Response:} Post-shot YOLO detection reports 5 visible balls compared to 6 in the pre-shot cache. The Hungarian algorithm identifies the unmatched cue ball as \textbf{Potted}. The restoration UI prompts ``Cue Ball in-pocket; restore to pre-shot coordinate'', directing the laser to the original strike spot.
\end{itemize}

\subsection{Test Case 4: Clean Legal Shot}
\begin{itemize}
    \item \textbf{Scenario:} Striker executes a legal safety shot. No foul is committed.
    \item \textbf{System Response:} Motion ceases, settlement confirms stability. With no foul invoked by the referee, the system automatically commits the post-shot resting state as the new \textit{Pre-Shot Baseline} for the incoming player, resetting the motion sensor for the subsequent stroke.
\end{itemize}

\section{Physical Rig and Web Interface Verification}
Comprehensive physical testing on the constructed table rig validated all engineering aspects:
\begin{itemize}
    \item The overhead PVC gantry maintained structural rigidity with less than $0.5\,\text{mm}$ deflection under ambient air currents.
    \item Controlled indoor ambient illumination combined with YOLOv8s learned chromatic features prevented false circle contours from cushion shadows.
    \item The Three.js 3D virtual table twin faithfully mirrored physical table states in real-time, providing an intuitive, interactive supervisory interface.
    \item The live servo nudge feature ($\pm 0.2^\circ$ fine-tuning) provided the referee with instant micro-trim adjustment capability, ensuring complete user trust and operational transparency.
\end{itemize}
